\documentclass[10pt,a4paper]{amsart}
\usepackage[margin=19mm]{geometry}
\usepackage{amsmath,amssymb,mathtools}
\usepackage[scr=rsfs,cal=euler]{mathalfa}
\usepackage{xfrac}
\usepackage[unicode,hidelinks,bookmarks=false]{hyperref}
\newcommand{\ie}{i.e.\ }
\newcommand{\cf}{cf.\ }
\newcommand{\resp}{respectively\ }
\newcommand{\Subsection}{\subsection}
\providecommand{\nobreakdash}{\nobreak\hbox{-}}
\setlength{\emergencystretch}{2em}
\multlinegap0pt
\usepackage{bm}
\usepackage{amssymb}
\usepackage{tikz}
\newtheorem{lem}{Lemma}
\usepackage{cleveref}
\crefname{lem}{Lemma}{Lemmas}
\DeclareMathOperator{\dr}{dr}
\title{Effective Velocities in the Toda Lattice}
\author{Amol Aggarwal}
\date{Source: arXiv:2503.11407v4 (CC BY 4.0)}
\begin{document}
\maketitle

\noindent\emph{Source and licence.} Effective Velocities in the Toda Lattice. Version arXiv:2503.11407v4 (CC BY 4.0), distributed by arXiv under the Creative Commons Attribution 4.0 International licence, http://creativecommons.org/licenses/by/4.0/, as recorded in the arXiv licence metadata for this exact version and shown on the abstract page https://arxiv.org/abs/2503.11407v4. Full licence notice: \texttt{licenses/arxiv-2503.11407-CC-BY-4.0.txt}.\par\smallskip

\noindent\textit{Editorial scope.} Counted unit: the article's lem derivativefg (source0.tex lines 724--731). The article's complete original proof of it is delivered with it (source0.tex lines 3237--3281, one \texttt{proof} environment of 45 lines, which the article places away from the statement). No mathematics is written, repaired or re-derived: the statement and the proof are the article's own lines, in the article's own notation, and no retained line is substituted or re-flowed.  Retained uncounted as the source-local context the proof needs: lines 269--273; lines 649--661; lines 665--670; lines 675--680; lines 713--721.  Nothing was omitted from any retained span, so the package carries no omission record.  No retained line contains a citation, so the wrapper carries no bibliography and no bibliography entry.  No retained line contains a figure environment, an image-inclusion command or drawing code, so no image is delivered and the package declares no asset dependency.  Editorial formatting only, in addition: an AMS \texttt{amsart} preamble that restates the article's own declarations and macros used by the retained lines, namely the environments \texttt{lem} on separate counters, together with the standard packages those lines need.  The retained environments print in this document as Lemma 1 in order of appearance, whereas the article prints the same items with the numbering of its own full document.  The complete native source of the article as distributed in the arXiv e-print for this exact version is bundled unchanged as \texttt{sources/174762/source0.tex}.
		\noindent For any function $h \in \mathcal{H}$, denote the associated multiplication operator by $\bm{h}$, defined by setting $\bm{h} f = hf$ for any $f : \mathbb{R} \rightarrow \mathbb{R}$; if $h \in \mathcal{H}$ is constant (that is, of the form $h = a \varsigma_0$ for some $a \in \mathbb{R}$), we identify $h = \bm{h}$. Further define the integral operator $\bm{\mathrm{T}}$ by setting  
		\begin{flalign}
			\label{operatort}
			\bm{\mathrm{T}} f (x) = 2 \displaystyle\int_{-\infty}^{\infty} \log |x-y| f(y) dy,
		\end{flalign}

	\begin{lem}
		
		\label{rhoexponential} 
		
		There exists a constant $C > 1$ such that 
		\begin{flalign}
			\label{estimate0} 
			\begin{aligned} 
			\varrho_{\beta} (x) \le C (|x &|+1)^{2 \theta} e^{-\beta x^2/2}; \qquad \varrho_{\beta}' (x) \le C(|x|+1) \varrho_{\beta} (x); \\
			& \varrho (x) \le C (|x|+1)^{2\theta+1} e^{-\beta x^2/2}.
			\end{aligned} 
		\end{flalign} 
	\end{lem} 

	\begin{lem}
		\label{rhorho} 
		
		For any $x \in \mathbb{R}$, we have $\varrho(x) = \theta \cdot ( \bm{\mathrm{T}} \varrho (x) + \alpha ) \cdot \varrho_{\beta} (x)$.
		
	\end{lem} 

	\begin{lem}
		\label{alphat}
		
		There exists a constant $c > 0$ such that $\bm{\mathrm{T}} \varrho (x) + \alpha > c$, for any $x \in \mathbb{R}$.
		
	\end{lem}

	\begin{lem} 
		\label{xf2} 
		
		There exists a constant $C>1$ such that the following holds. For any function $f \in \mathcal{H}$ and real number $x \in \mathbb{R}$, we have 
		\begin{flalign*} 
			|f^{\dr} (x)| \le C \cdot |f(x)| + C \| f \|_{\mathcal{H}} \cdot \log (|x|+2).
		\end{flalign*} 
		
	\end{lem} 

	\begin{lem} 
		\label{derivativefg}
		
		There exists a constant $C>1$ such that the following holds. For any differentiable function $f \in \mathcal{H}$ such that $f' \in \mathcal{H}$, and real number $x \in \mathbb{R}$, we have 
		\begin{flalign*}
			|\partial_x f^{\dr} (x)| \le C \cdot |f'(x)| + C (  \| f' \|_{\mathcal{H}} + \| f \|_{\mathcal{H}} ) \cdot \log (|x|+2).
		\end{flalign*} 
	\end{lem} 

	\begin{proof}[Proof of \Cref{derivativefg}]
		
		Denote $g(x) = f^{\dr} (x)$ for each $x \in \mathbb{R}$. Then, observe that 
		\begin{flalign}
			\label{gderivativef}
			g' = (f' + \bm{\mathrm{T} \varrho_{\beta}'} g)^{\dr}.
		\end{flalign} 
		
		\noindent Indeed, differentiating the relation $g(x) = \theta f(x) + \theta \cdot \bm{\mathrm{T} \varrho_{\beta}} g (x)$ in $x$ (and recalling \eqref{operatort}) yields 
		\begin{flalign*}
			g'(x) = \theta f'(x) + 2 \theta \displaystyle\int_{-\infty}^{\infty} \log |y| \cdot \partial_x \big( g(x-y) \varrho_{\beta} (x-y) \big) dy = \theta f' + \theta \bm{\mathrm{T} \varrho_{\beta}} g' + \theta \bm{\mathrm{T} \varrho_{\beta}'} g,
		\end{flalign*} 
		
		\noindent from which we obtain \eqref{gderivativef}. Applying \Cref{xf2}, it follows there exists a constant $C_1>1$ such that, for any $x \in \mathbb{R}$, 
		\begin{flalign}
			\label{gderivativexestimate0} 
			|g'(x)| \le C_1 \cdot ( |f'(x)| + |\bm{\mathrm{T} \varrho_{\beta}'} g(x)| ) + C_1 \cdot ( \| f' \|_{\mathcal{H}} + \| \bm{\mathrm{T} \varrho_{\beta}'} g \|_{\mathcal{H}} ) \cdot \log (|x|+2).
		\end{flalign}
		
		We must now bound $\bm{\mathrm{T} \varrho_{\beta}'} g$. To that end, observe from \eqref{operatort} that 
		\begin{flalign}
			\label{tderivativeg} 
			\begin{aligned} 
			|\bm{\mathrm{T} \varrho_{\beta}'} g(x)| & \le 2  \displaystyle\int_{-\infty}^{\infty} \log |x-y| \cdot \varrho_{\beta}' (y) |g(y)| dy  \\
			& \le 2 \Bigg( \displaystyle\int_{-\infty}^{\infty} |g(y)|^2 \varrho (y) dy \Bigg)^{1/2} \Bigg( \displaystyle\int_{-\infty}^{\infty} (\log |x-y|)^2 \varrho (y)^{-1} \varrho_{\beta}' (y)^2 dy \Bigg)^{1/2}.
			\end{aligned}
		\end{flalign} 
		
		\noindent From \Cref{rhoexponential}, \Cref{rhorho}, and \Cref{alphat}, there exists constants $C_2 > 1$ and $C_3 >1$ that 
		\begin{flalign*}
			|\varrho_{\beta}' (y)| \le C_2 (|y|+1) \cdot \varrho_{\beta} (y) \le C_3 (|y|+1) \cdot \varrho (y).
		\end{flalign*}
		
		\noindent Inserting this into \eqref{tderivativeg} (and using the boundedness and exponential decay of $\varrho$, from \Cref{rhoexponential}) yields for some constants $C_4 > 1$ and $C_5 > 1$ that
		\begin{flalign*} 
			|\bm{\mathrm{T} \varrho_{\beta}'} g(x)| \le C_4 \| g \|_{\mathcal{H}} \cdot \Bigg( \displaystyle\int_{-\infty}^{\infty} (\log |x-y|)^2 (|y|+1) \varrho (y) dy \Bigg)^{1/2} \le C_5 \| g \|_{\mathcal{H}} \cdot \log (|x|+2).
		\end{flalign*} 
		
		\noindent By \Cref{xf2}, we have $\| g \|_{\mathcal{H}} \le C_6 \| f \|_{\mathcal{H}}$ for some $C_6 > 1$, so it follows for some $C_7 > 1$ that 
		\begin{flalign*}
			|\bm{\mathrm{T} \varrho_{\beta}'} g(x)| \le C_7 \| f \|_{\mathcal{H}} \cdot \log (|x|+2); \qquad	\| \bm{\mathrm{T} \varrho_{\beta}'} g \|_{\mathcal{H}} \le C_7 \| f \|_{\mathcal{H}}.
		\end{flalign*}
		
		\noindent Applying these bounds in \eqref{gderivativexestimate0} then yields the lemma.
	\end{proof} 

\end{document}
